% !TEX root = ../../tir_monograph_v12.tex
\chapter{Euclidean Geometry and the Spatial Dimension Gate}
\label{ch:v12_euclidean_dimension}

\section{Canonical affine relation carrier}

The normalized two-level state carrier has affine hull
\[
\mathcal A_2
=
\{\rho=\rho^\dagger:\operatorname{Tr}\rho=1\}
=
\frac12I+\operatorname{Herm}_0(2).
\]
Write
\[
V:=\operatorname{Herm}_0(2)
=
\operatorname{span}_{\mathbb R}\{\sigma_x,\sigma_y,\sigma_z\}.
\]
Since \(\dim_{\mathbb R}\operatorname{Herm}(2)=4\) and the trace-one condition removes one
real affine degree of freedom,
\[
\boxed{\dim_{\mathbb R}V=3.}
\]

For an ordered pair of admitted point states \((\rho_x,\rho_y)\), the affine torsor
has a unique endpoint-carrying displacement
\[
\boxed{\mathfrak e_{xy}:=\rho_y-\rho_x\in V},
\qquad
\rho_x+\mathfrak e_{xy}=\rho_y.
\]
The generator-normalized relation used below is
\[
\boxed{\mathcal E_{xy}:=2\mathfrak e_{xy}.}
\]
It satisfies the exact endpoint identities
\[
\mathcal E_{yx}=-\mathcal E_{xy},
\qquad
\boxed{\mathcal E_{xz}=\mathcal E_{xy}+\mathcal E_{yz}},
\]
and therefore
\[
\mathcal E_{xy}+\mathcal E_{yz}+\mathcal E_{zx}=0.
\]
This endpoint-carrying universal property is the canonical local relation exported
by the v12 spatial branch.

\section{Invariant Euclidean metric}

On \(V\), define
\[
\boxed{
\langle A,B\rangle
=
\frac12\operatorname{Tr}(AB).
}
\]
The Pauli basis obeys
\[
\frac12\operatorname{Tr}(\sigma_i\sigma_j)=\delta_{ij}.
\]
Hence, for
\[
A=\mathbf a\cdot\boldsymbol\sigma,
\qquad
B=\mathbf b\cdot\boldsymbol\sigma,
\]
one has
\[
\boxed{
\langle A,B\rangle=\mathbf a\cdot\mathbf b,
\qquad
\|A\|^2=\frac12\operatorname{Tr}(A^2)=|\mathbf a|^2.
}
\]
The positive quadratic form is therefore Euclidean in Pauli coefficients.

The Pythagorean relation follows immediately. If
\[
\langle A,B\rangle=0,
\]
then
\[
\begin{aligned}
\|A+B\|^2
&=\langle A+B,A+B\rangle\\
&=\|A\|^2+2\langle A,B\rangle+\|B\|^2,
\end{aligned}
\]
so
\[
\boxed{\|A+B\|^2=\|A\|^2+\|B\|^2.}
\]
Thus the local affine relation carrier already contains the Euclidean norm,
orthogonality and Pythagorean closure used by the later spatial construction.

\section{Rotational covariance}

For \(U\in SU(2)\),
\[
A\mapsto UAU^\dagger
\]
preserves trace and the Hilbert--Schmidt metric.  The induced action on coefficient
vectors is
\[
\boxed{
\operatorname{Ad}:SU(2)\longrightarrow SO(3),
\qquad
SU(2)/\{\pm I\}\cong SO(3).
}
\]
The unit relation directions therefore form
\[
\boxed{
\{A\in V:\|A\|=1\}\cong S^2.
}
\]

\section{Rank-three stabilization}

At a relational locus \(x\), define the outgoing relation span
\[
W_x
:=
\operatorname{span}_{\mathbb R}
\{\mathcal E_{xy}:y\sim x\}
\subseteq V.
\]
Assume a populated primitive patch,
\[
W_x\neq\{0\},
\]
and full local relational isotropy,
\[
R(W_x)=W_x
\qquad
\text{for every }R\in SO(3)
\]
from the adjoint action above.

The defining real representation of \(SO(3)\) on \(\mathbb R^3\) is irreducible.
Its invariant linear subspaces are therefore \(\{0\}\) and the full carrier.
Consequently
\[
\boxed{
W_x\neq0,\quad SO(3)W_x=W_x
\Longrightarrow
W_x=V
\Longrightarrow
\dim W_x=3.
}
\]
Rank one selects an axis and rank two selects a plane; the full primitive isotropy
stabilizes the three-dimensional carrier.

\section{Minimal faithful realization crosscheck}

The same dimension is independently fixed by the minimal faithful real orthogonal
realization of the primitive rotation group. A continuous \(SO(3)\) action on a
one-dimensional real orthogonal carrier has trivial connected image.  A
two-dimensional orthogonal image lies in \(SO(2)\), while \(SO(3)\) is non-Abelian
and perfect.  The defining three-dimensional representation is faithful.

Therefore the minimal faithful real orthogonal carrier has
\[
\boxed{\dim_{\mathbb R}V_x=3}
\]
and is unique up to orthogonal equivalence.  Its invariant positive metric is
unique up to an overall positive scale.  The canonical dimensionless
normalization used here is exactly
\[
h(A,B)=\frac12\operatorname{Tr}(AB).
\]

\section{Continuum interface}

The discrete spatial branch is now source-bound by
\[
\boxed{
(\rho_x,\rho_y)
\longmapsto
\mathcal E_{xy}=2(\rho_y-\rho_x)
\in\operatorname{Herm}_0(2)\cong\mathbb R^3.
}
\]
The remaining geometric completion gate is the regular gluing/refinement step
that turns neighboring local copies of \(V\) into a smooth coframe/tangent-bundle
carrier.  Chapter~\ref{ch:v12_connections_holonomy} supplies the discrete
connection, solder and torsion objects required for that refinement.

\section{Source and validation receipt}

The canonical theorem surfaces are
\begin{itemize}
\item \path{TIR/subrepos/the-space-of-geometry/foundations/CANONICAL_SPATIAL_RELATION_EXTRACTION_V0_1.md};
\item \path{TIR/foundations/TIR_RANK3_ISOTROPY_STABILIZATION_V0_1.md};
\item \path{TIR/subrepos/the-space-of-geometry/foundations/SPATIAL_PROMOTION_UNIQUENESS_V0_1.md}.
\end{itemize}

The deterministic v12 receipt is generated by
\[
\path{TIR/validation/tir_v12_ch05_euclidean_spatial_gate_v0_1.py}.
\]
